% !TeX root = ../教材题库_全部条目_转写工作稿.tex
% Source: XJ-B1-C06-HW-P232-Y0260
% 第6章 / 6.3 统计图表 / 教材 p.232 / PDF 第239页
% OCR未校对：公式与文字均可能有误；未完成数学TeX转写；禁止直接发给学生
% 原题及相邻原解截图（核对用）：../题目原图/XJ-B1-C06-HW-P232-Y0260-01.webp
% 本文件只在转写完成后作为题干使用；OCR 条目尚不是可靠题干。
\begin{bankitem}{习题6.3 第4题}{第6章 / 6.3 统计图表 / 教材 p.232 / PDF 第239页}{XJ-B1-C06-HW-P232-Y0260}{OCR未校对：公式与文字均可能有误；未完成数学TeX转写；禁止直接发给学生}
\begin{ocrdraft}
[教材p.232 / PDF第239页]
\par
4．如图是根据我国部分城市某年6月份的平均气温数据得到的样本频率分布直
\par
方图，其中平均气温的范围是[20，26]，样本数据的分组为[20，(1)，[21，22），
\par
[22，23），[23，24），[24，25），[25，26]．已知样本中平均气温低于22\ensuremath{{}^\circ\mathrm C}的城
\par
市个数为11，求样本中平均气温不低于25\ensuremath{{}^\circ\mathrm C}的城市个数．
\par
组距
\par
0．26
\par
0、22
\par
0．18
\par
0．12
\par
0．10
\par
0
\par
20212223242526平均气温／\ensuremath{{}^\circ\mathrm C}
\par
（第4题）
\end{ocrdraft}
\end{bankitem}
